\scanpage{14}
\mspage{- 2 -}\hfill n\handwritten{°354}

duation de type $\handwritten{\Delta}'$ sur $G$.

Par exemple, soient $\handwritten{\Delta}_{\unsure{o|0}}$ un monoïde commutatif ayant un élément neutre 0,
$I$ un ensemble d'indices arbitraire, $\handwritten{\Delta}=\handwritten{\Delta}_o^{(I)}$, sous-monoïde du produit
$\handwritten{\Delta}_0^I$ formé des familles $(\handwritten{\lambda}_{\handwritten{\unsure{\iota|i}}})_{\handwritten{\iota}\handwritten{\in}I}$ à support fini. Soit $(n_{\handwritten{\iota}})_{\handwritten{\iota}\handwritten{\in}I}$ une famille
quelconque d'entiers $\handwritten{\geq}0$ ; l'application $(\handwritten{\lambda}_{\handwritten{\iota}})\handwritten{\to}\handwritten{\sum}\limits_{\handwritten{\iota}\handwritten{\in}I}n_{\handwritten{\iota}}\handwritten{\lambda}_{\handwritten{\iota}}$ est un homomorphisme $\handwritten{\rho}$ de $\handwritten{\Delta}$ dans $\handwritten{\Delta}_o$ tel que $\handwritten{\rho}(0)=0$, et associe donc à toute graduation de type $\handwritten{\Delta}$ une graduation de type $\handwritten{\Delta}_o$.

3) Plus particulièrement, si dans l'exemple 2), on prend $\handwritten{\Delta}=\handwritten{\Delta}_o\handwritten{\times}\handwritten{\Delta}_o$, on
dit qu'une graduation de $\handwritten{\Delta}$ est une \uline{bigraduation} de type $\handwritten{\Delta}_o$ ; les trois
homomorphismes $(\handwritten{\lambda},\mu)\handwritten{\to}\handwritten{\lambda}$, $(\handwritten{\lambda},\mu)\handwritten{\to}\mu$, $(\handwritten{\lambda},\mu)\handwritten{\to}\handwritten{\lambda}+\mu$ associent alors
à toute bigraduation de type $\handwritten{\Delta}_o$ trois graduations de type $\handwritten{\Delta}_o$ dites respectivement \uline{première} et \uline{seconde graduation partielle}, et \uline{graduation totale}. De
façon plus explicite, si $(G_{\handwritten{\lambda}\mu})_{(\handwritten{\lambda},\mu)\handwritten{\in}\handwritten{\Delta}}$ est la bigraduation donnée, la
première graduation partielle est formée des $G'_{\handwritten{\lambda}}=\handwritten{\bigoplus}\limits_{\mu\handwritten{\in}\handwritten{\Delta}_o}G_{\handwritten{\lambda}\mu}$, la seconde
des $G''_{\handwritten{\lambda}}=\handwritten{\bigoplus}\limits_{\mu\handwritten{\in}\handwritten{\Delta}_o}G_{\mu \handwritten{\lambda}}$ et la graduation totale des $G'''_{\handwritten{\lambda}}=\handwritten{\bigoplus}\limits_{\mu+\unsure{y|\nu}=\handwritten{\lambda}}G_{\mu y}$. On
notera que l'on a $G'_{\handwritten{\lambda}}\handwritten{\cap}G''_\mu=G_{\handwritten{\lambda}\mu}$ ; inversement, si $(G'_{\handwritten{\lambda}})$ et $(G''_{\handwritten{\lambda}})$ sont deux
graduations de type $\handwritten{\Delta}_o$ sur un groupe commutatif $G$, telles que $G$ soit somme
directe des sous-groupes $G'_{\handwritten{\lambda}}\handwritten{\cap}G''_\mu=G_{\handwritten{\lambda}\mu}$, ces sous-groupes définissent une
bigraduation de type $\handwritten{\Delta}$ sur $G$, dont $(G'_{\handwritten{\lambda}})$ et $(G''_{\handwritten{\lambda}})$ sont les deux graduations
partielles.

Nous laissons au lecteur le soin de généraliser ces définitions au cas o\unsure{\handwritten{\`}{u}|\handwritten{\=}{u}}
$\handwritten{\Delta}=\handwritten{\Delta}_o^n$ .

DEFINITION 2.- \uline{Etant donnés un anneau} $A$ \uline{et une graduation} $(A_{\handwritten{\lambda}})$ \uline{de type} $\handwritten{\Delta}$
\uline{sur le groupe additif} $A$, \uline{on dit que cette graduation est compatible avec la}
\uline{structure d'anneau de} $A$ \uline{si l'on a}
\begingroup
\makeatletter\tagsleft@true\makeatother
\[
 A_{\handwritten{\lambda}}A_\mu\handwritten{\subset}A_{\handwritten{\lambda}+\mu}\qquad
 \text{\uline{quels que soient} }\handwritten{\lambda}\ \text{\uline{et} }\mu\ \text{\uline{dans} }\handwritten{\Delta}\ .
 \tag*{(1)}
\]
\endgroup
% archive sheet number: 13
